\documentclass[12pt,a4paper]{article}

% === Packages ===
\usepackage[utf8]{inputenc}
\usepackage[T1]{fontenc}
\usepackage{lmodern}
\usepackage[margin=1in]{geometry}
\usepackage{amsmath,amssymb,amsfonts}
\usepackage{graphicx}
\usepackage{booktabs}
\usepackage{longtable}
\usepackage{multirow}
\usepackage{array}
\usepackage{tabularx}
\usepackage{hyperref}
\usepackage{cleveref}
\usepackage{natbib}
\usepackage{setspace}
\usepackage{enumitem}
\usepackage{xcolor}
\usepackage{subcaption}
\usepackage{float}

\hypersetup{
    colorlinks=true,
    linkcolor=blue!70!black,
    citecolor=green!50!black,
    urlcolor=blue!70!black
}

\onehalfspacing

% === Title ===
\title{Machine Learning in Quantitative Finance:\\A Systematic Review of Methods, Applications, and Open Challenges (2015--2025)}

\author{
    Akram Khan\thanks{Independent Researcher. ORCID: \href{https://orcid.org/0009-0002-7521-8648}{0009-0002-7521-8648}. Email: akramkhan@ucla.edu}
}

\date{\today}

\begin{document}

\maketitle

\begin{abstract}
We review 227 studies applying machine learning to quantitative finance, published between 2015 and 2025, and organize them into a two-dimensional taxonomy covering six application domains (return prediction, risk management, portfolio optimization, derivatives pricing, market microstructure, and alternative data) and five methodological families (tree-based methods, deep learning, reinforcement learning, NLP, and hybrids). The literature reveals three generational shifts: gradient-boosted trees establishing dominance on tabular financial data (2015--2018), attention mechanisms and transformers displacing recurrent networks (2018--2022), and large language models entering financial applications (2022--2025). Our main findings are that tree-based ensembles remain hard to beat for cross-sectional prediction, that the strongest ML improvements over traditional methods are in derivatives pricing and volatility forecasting rather than return prediction, that reinforcement learning for portfolio management has attracted more papers than evidence, and that only about one in five studies shares code. We provide standardized comparison tables across all domains and conclude with ten open problems, arguing that the field's most pressing need is not better models but better benchmarks and reproducibility standards.

\medskip
\noindent\textbf{Keywords:} machine learning, quantitative finance, systematic review, deep learning, return prediction, portfolio optimization, risk management, NLP in finance

\medskip
\noindent\textbf{JEL Classification:} G11, G12, G14, G17, C45, C55
\end{abstract}

\newpage
\tableofcontents
\newpage

% === Section 1: Introduction ===
\section{Introduction}\label{sec:introduction}

Ten years ago, applying machine learning to financial markets was a niche
pursuit. A handful of papers tested random forests or SVMs on stock
returns, usually on one market and one horizon, and the results were mixed
enough that most finance departments shrugged them off. That is no longer
the case. By 2025, ML has reached every corner of quantitative finance.
Cross-sectional asset pricing. Options hedging. Earnings-call sentiment.
The volume of published work has grown to the point where keeping track of
it is itself a research problem.

The trouble is that the work is scattered. An ML researcher looking for
the best model to predict equity returns will find relevant papers at
NeurIPS, in the \textit{Review of Financial Studies}, on arXiv's q-fin
section, and buried in SSRN working-paper series. The conventions of each
venue are different. ML papers report out-of-sample $R^2$ and treat a
0.5\% improvement as a win. Finance papers want Sharpe ratios net of
transaction costs and a comparison against Fama-French factor models. A
result that looks impressive by one community's standards may look
unremarkable, or even suspicious, by the other's. The fragmentation wastes
effort: researchers in one camp routinely rediscover what the other camp
settled years ago.

\subsection{What This Survey Does}

This paper tries to impose some order on the chaos. The survey covers ML
applications in quantitative finance from 2015 through 2025: 227 studies
across six application domains (return prediction, risk management,
portfolio optimisation, derivatives pricing, market microstructure, and
alternative data/NLP). Each study is classified along two axes,
methodology and application, which produces a taxonomy matrix that shows
at a glance where the literature is dense and where it has holes.

Three things distinguish this survey from its predecessors.

\begin{enumerate}[leftmargin=*]
    \item \textbf{Standardised comparison tables.} For each domain we report the dataset, sample period, ML method, benchmark, performance metric, and, critically, whether code and data are public. Most prior surveys describe methods in prose. We tabulate them so a reader can compare directly.

    \item \textbf{Full coverage of the LLM era.} Existing surveys largely predate the 2022--2025 explosion of large language models in finance. We cover BloombergGPT, FinBERT, GPT-based return prediction, and the broader question of whether foundation models will reshape the field.

    \item \textbf{A reproducibility audit.} For each study we track whether code and data are available. The numbers are not encouraging. Of the 227 studies in our sample, roughly 47 (21\%) share replication code; on the 119-study subset released as a structured CSV alongside this paper the count is 33 of 119 (28\%). The gap is a sample selection effect: the released CSV is enriched for studies whose code-availability flag we have personally verified through GitHub or repository inspection. We flag this not as an aside but as a first-order problem for the field.
\end{enumerate}

\subsection{Prior Surveys and Where They Fall Short}

The literature already has several surveys, each useful but each with blind spots. \citet{Henrique2019} cover ML for stock prediction through 2018, but that cutoff misses transformers, LLMs, and the post-pandemic data regime entirely. \citet{Ozbayoglu2020} survey deep learning architectures for finance, though their emphasis is on describing network structures rather than comparing empirical results head-to-head. \citet{DePrado2018} and \citet{Dixon2020} wrote valuable textbooks, but textbooks age, and neither could anticipate the pace of change after 2020.

Among more recent work, \citet{KellyXiu2023} provide the most authoritative treatment; their \textit{Foundations and Trends in Finance} monograph is close to definitive on return prediction and factor models. But their scope is deliberately narrow: they do not cover NLP, derivatives pricing, or microstructure. \citet{NazarethReddy2023} cast a wider net (126 articles, 44 journals) but focus on categorization rather than critical evaluation, telling you what methods have been applied where rather than which ones actually work. \citet{Goodell2022} take a bibliometric approach (348 papers, co-citation analysis) that maps the intellectual landscape but provides no guidance to a practitioner wondering which model to deploy.

Our survey fills the gap between these approaches: broad enough to cover all six domains, deep enough to include comparison tables and reproducibility tracking, and current enough to address the LLM wave.

\subsection{Roadmap}

\Cref{sec:methodology} describes our search strategy and inclusion criteria. \Cref{sec:taxonomy} lays out the two-dimensional taxonomy. \Cref{sec:return_prediction} through \Cref{sec:alternative_data} work through each application domain in turn. \Cref{sec:cross_cutting} tackles four themes that cut across domains: interpretability, non-stationarity, compute, and the reproducibility crisis. \Cref{sec:open_problems} identifies ten open problems. \Cref{sec:conclusion} steps back and offers an overall assessment.


% === Section 2: Methodology ===
\section{Systematic Review Methodology}\label{sec:methodology}

We follow the PRISMA (Preferred Reporting Items for Systematic Reviews and Meta-Analyses) guidelines adapted for computational research \citep{Page2021}. This section details our search strategy, selection criteria, and data extraction protocol.

\subsection{Search Strategy}

We conducted systematic searches across five databases:
\begin{itemize}[leftmargin=*]
    \item \textbf{Scopus} --- comprehensive coverage of finance and CS journals
    \item \textbf{Web of Science} --- primary citation database for impact tracking
    \item \textbf{SSRN} --- working papers and preprints in finance
    \item \textbf{arXiv (q-fin, cs.LG, stat.ML)} --- preprints in quantitative finance and ML
    \item \textbf{Google Scholar} --- supplementary search for grey literature and conference proceedings
\end{itemize}

Search queries combined financial terms (\texttt{``stock prediction'' OR ``portfolio optimization'' OR ``risk management'' OR ``options pricing'' OR ``market microstructure'' OR ``quantitative finance''}) with ML terms (\texttt{``machine learning'' OR ``deep learning'' OR ``neural network'' OR ``random forest'' OR ``gradient boosting'' OR ``reinforcement learning'' OR ``natural language processing'' OR ``transformer'' OR ``LSTM''}).

Initial searches were conducted in March 2026, with supplementary forward-citation searches completed in April 2026.

\subsection{Inclusion and Exclusion Criteria}

\begin{table}[H]
\centering
\caption{Inclusion and exclusion criteria}
\label{tab:criteria}
\begin{tabularx}{\textwidth}{lX}
\toprule
\textbf{Include} & \textbf{Exclude} \\
\midrule
Published 2015--2025 & Published before 2015 \\
Applies ML to a quantitative finance problem & Pure ML methodology with no financial application \\
Presents empirical results on financial data & Theoretical/conceptual only \\
Journal article, conference paper, or working paper & Theses, textbooks, blog posts \\
English language & Non-English \\
Uses real or realistic simulated financial data & Uses only toy/synthetic examples \\
\bottomrule
\end{tabularx}
\end{table}

\subsection{Selection Process}

\begin{enumerate}[leftmargin=*]
    \item \textbf{Identification:} Database searches returned 4,783 unique records after deduplication. The arXiv q-fin section alone contributed over 12,600 submissions during our review period (2015--2025), growing from 852 total submissions in 2015 to 2,435 in 2024---a nearly threefold increase reflecting the rapid growth of quantitative finance research \citep{arXivStats2025}. Scopus, Web of Science, and SSRN contributed the remainder.
    \item \textbf{Screening:} Title and abstract screening against inclusion criteria reduced the set to 847 papers.
    \item \textbf{Eligibility:} Full-text review of screened papers identified 196 eligible studies.
    \item \textbf{Inclusion:} Forward and backward citation searches added 31 additional papers, yielding a final sample of 227 studies.
\end{enumerate}

\subsection{Data Extraction Protocol}

For each included study, we extracted the following fields into a structured database:

\begin{table}[H]
\centering
\caption{Data extraction fields}
\label{tab:extraction}
\begin{tabularx}{\textwidth}{lX}
\toprule
\textbf{Field} & \textbf{Description} \\
\midrule
Authors \& Year & Standard citation information \\
Venue & Journal, conference, or preprint server \\
Application Domain & Return prediction, risk management, portfolio optimization, derivatives, microstructure, or alternative data \\
ML Method(s) & Primary and benchmark methods used \\
Asset Class & Equities, fixed income, FX, commodities, crypto, multi-asset \\
Data Source & Specific database or data provider \\
Sample Period & Start and end dates of empirical analysis \\
Geography & US, Europe, Asia, emerging markets, global \\
Primary Metric & Main performance measure reported \\
Key Finding & One-sentence summary of the main result \\
Code Available & Yes/No with URL if available \\
Data Available & Public / Proprietary / Mixed \\
\bottomrule
\end{tabularx}
\end{table}

\subsection{Quality Assessment}

We do not exclude studies based on quality scores, as methodological standards differ substantially between the ML and finance communities. Instead, we flag quality dimensions that readers should consider when interpreting results:

\begin{itemize}[leftmargin=*]
    \item \textbf{Out-of-sample testing:} Does the study use a proper temporal train/validation/test split, or does it employ cross-validation (which can introduce look-ahead bias in time series)?
    \item \textbf{Transaction costs:} Are trading strategies evaluated net of realistic transaction costs?
    \item \textbf{Benchmark comparison:} Is the ML model compared against appropriate traditional benchmarks (e.g., GARCH for volatility, Fama-French for returns)?
    \item \textbf{Statistical significance:} Are results tested for statistical significance, or are only point estimates reported?
    \item \textbf{Reproducibility:} Is code and/or data publicly available?
\end{itemize}

These dimensions are reported in our comparison tables (\Cref{sec:appendix_tables}) to help readers assess the strength of evidence for each finding.


% === Section 3: Taxonomy of Methods ===
\section{A Taxonomy of Machine Learning in Finance}\label{sec:taxonomy}

We organize the literature along two dimensions: \textit{methodology} (the ML technique) and \textit{application domain} (the financial problem). This section describes each dimension and presents the resulting taxonomy matrix.

\subsection{Methodological Classification}

We classify ML methods into five families, reflecting the major paradigms that have been applied in quantitative finance.

\subsubsection{Tree-Based Ensemble Methods}

Tree-based ensembles have emerged as the dominant methodology for structured (tabular) financial data. The family includes random forests \citep{Breiman2001}, gradient-boosted decision trees \citep{Friedman2001}, and their modern implementations: XGBoost \citep{Chen2016}, LightGBM \citep{Ke2017}, and CatBoost \citep{Prokhorenkova2018}. Their advantages in financial applications include:

\begin{itemize}[leftmargin=*]
    \item Natural handling of heterogeneous features (continuous, categorical, ordinal)
    \item Built-in feature importance measures that support interpretability
    \item Robustness to outliers and non-linear relationships
    \item Competitive performance without extensive hyperparameter tuning
    \item Ability to capture interaction effects without explicit feature engineering
\end{itemize}

The empirical dominance of GBDT in tabular data settings has been confirmed in large-scale benchmarks \citep{Grinsztajn2022, ShwartzZiv2022alt}, and financial applications consistently reflect this finding. In cross-sectional return prediction, GBDT methods match or exceed deep learning alternatives in the majority of studies reviewed (\Cref{sec:return_prediction}).

\subsubsection{Deep Learning Architectures}

Deep learning in finance has progressed through three distinct generations:

\textbf{First generation (2015--2018): Recurrent networks.} Long Short-Term Memory (LSTM) networks \citep{Hochreiter1997} and Gated Recurrent Units (GRU) \citep{Cho2014} were the first deep learning architectures widely applied to financial time series. Their ability to model sequential dependencies made them natural candidates for price prediction and volatility forecasting.

\textbf{Second generation (2018--2022): Attention and transformers.} The transformer architecture \citep{Vaswani2017} and its variants have progressively displaced recurrent models. Financial applications include temporal fusion transformers \citep{Lim2021}, informer architectures for long-horizon forecasting \citep{Zhou2021}, and custom attention mechanisms for cross-sectional asset pricing \citep{Chen2024}.

\textbf{Third generation (2022--2025): Foundation models.} Pre-trained large language models (GPT, LLaMA, and domain-specific models such as FinBERT \citep{Araci2019}, BloombergGPT \citep{Wu2023}) have opened new approaches to financial text analysis, event extraction, and even direct numerical forecasting through in-context learning.

Additional deep learning approaches applied in finance include convolutional neural networks (CNNs) for chart pattern recognition and volatility surface modeling, autoencoders for anomaly detection and factor extraction, and generative adversarial networks (GANs) for synthetic data generation and scenario simulation.

\subsubsection{Reinforcement Learning}

Reinforcement learning (RL) frames financial decision-making as a Markov Decision Process where an agent learns to take actions (trades) to maximize cumulative reward (risk-adjusted returns). Applications include:

\begin{itemize}[leftmargin=*]
    \item Portfolio allocation and rebalancing \citep{Jiang2017, Ye2020}
    \item Optimal execution and market making \citep{Ning2021}
    \item Options hedging \citep{Buehler2019, Cao2021}
\end{itemize}

Key RL frameworks applied in finance include Deep Q-Networks (DQN), Proximal Policy Optimization (PPO), Soft Actor-Critic (SAC), and model-based approaches. A persistent challenge is the sim-to-real gap: policies learned in backtesting environments often fail to transfer to live markets due to non-stationarity and market impact.

\subsubsection{Natural Language Processing and LLMs}

NLP methods extract information from unstructured financial text: earnings calls, SEC filings, news articles, social media, and analyst reports. The field has evolved from bag-of-words and sentiment dictionaries \citep{Loughran2011} through word embeddings (Word2Vec, GloVe) to transformer-based models:

\begin{itemize}[leftmargin=*]
    \item \textbf{FinBERT} \citep{Araci2019}: BERT fine-tuned on financial text for sentiment classification
    \item \textbf{BloombergGPT} \citep{Wu2023}: 50B-parameter model trained on Bloomberg's financial data corpus
    \item \textbf{General-purpose LLMs}: GPT-4, Claude, and LLaMA applied zero-shot or few-shot to financial tasks
\end{itemize}

\subsubsection{Hybrid and Ensemble Approaches}

Many recent studies combine multiple methodologies:
\begin{itemize}[leftmargin=*]
    \item LSTM + attention mechanisms for enhanced sequence modeling
    \item NLP sentiment features fed into GBDT return prediction models
    \item RL agents using deep learning for state representation
    \item Stacking ensembles combining tree-based and neural network predictions
\end{itemize}

\subsection{Application Domain Classification}

We identify six application domains that collectively span the quantitative finance literature:

\begin{enumerate}[leftmargin=*]
    \item \textbf{Return prediction} (\Cref{sec:return_prediction}): Forecasting asset returns at various horizons (intraday to monthly), both cross-sectional (which stocks will outperform?) and time-series (will the market go up?).

    \item \textbf{Risk management} (\Cref{sec:risk_management}): Volatility forecasting, Value-at-Risk estimation, Expected Shortfall, credit risk modeling, and systemic risk measurement.

    \item \textbf{Portfolio optimization} (\Cref{sec:portfolio_optimization}): ML-enhanced portfolio construction, including direct weight learning, covariance estimation, and transaction cost optimization.

    \item \textbf{Derivatives pricing and hedging} (\Cref{sec:derivatives}): Option pricing beyond Black-Scholes, implied volatility surface modeling, and dynamic hedging strategies.

    \item \textbf{Market microstructure} (\Cref{sec:microstructure}): Order flow prediction, optimal execution, market making, and limit order book modeling.

    \item \textbf{Alternative data and NLP} (\Cref{sec:alternative_data}): Extraction of tradeable signals from news, social media, satellite imagery, and other non-traditional data sources.
\end{enumerate}

\subsection{The Taxonomy Matrix}

\Cref{tab:taxonomy_matrix} presents the two-dimensional taxonomy, indicating the approximate number of studies in our sample at each intersection and our qualitative assessment of research maturity.

\begin{table}[H]
\centering
\caption{Taxonomy matrix: Methods $\times$ Applications. Cell entries indicate approximate study count and maturity level (H = High, M = Medium, L = Low). Studies that apply more than one methodology (e.g., LSTM with attention, NLP features feeding a tree ensemble) appear in multiple rows; for this reason the cell totals (344) exceed the unique study count (227).}
\label{tab:taxonomy_matrix}
\small
\begin{tabular}{l c c c c c c}
\toprule
& \textbf{Returns} & \textbf{Risk} & \textbf{Portfolio} & \textbf{Derivatives} & \textbf{Microstr.} & \textbf{Alt.\ Data} \\
\midrule
\textbf{Tree-based}  & 40+ (H) & 15 (M) & 10 (M) & 5 (L)  & 5 (L)  & 10 (M) \\
\textbf{Deep learning} & 35+ (H) & 20 (M) & 15 (M) & 15 (M) & 10 (M) & 20 (M) \\
\textbf{RL}           & 5 (L)   & 3 (L)  & 20 (M) & 10 (M) & 8 (M)  & 2 (L) \\
\textbf{NLP / LLM}   & 15 (M)  & 5 (L)  & 3 (L)  & 2 (L)  & 2 (L)  & 30+ (H) \\
\textbf{Hybrid}       & 10 (M)  & 5 (L)  & 8 (M)  & 5 (L)  & 3 (L)  & 8 (M) \\
\bottomrule
\end{tabular}
\end{table}

Several patterns emerge from this matrix:

\begin{itemize}[leftmargin=*]
    \item \textbf{Return prediction} is the most studied intersection, dominated by tree-based and deep learning methods.
    \item \textbf{RL for portfolio optimization} has attracted disproportionate attention relative to its empirical success, suggesting a gap between theoretical appeal and practical viability.
    \item \textbf{NLP/LLM methods beyond sentiment} (e.g., for risk management or derivatives) remain under-explored and represent opportunities for future work.
    \item \textbf{Market microstructure} is under-represented in the ML literature relative to its importance in practice, likely due to data access barriers.
    \item \textbf{Tree-based methods for derivatives} are surprisingly sparse, despite the tabular nature of options data.
\end{itemize}


% === Section 4: Return Prediction ===
\section{Return Prediction}\label{sec:return_prediction}

Return prediction is the oldest and most heavily studied application of ML in finance. We organize this section into cross-sectional prediction (which assets will outperform), time-series prediction (will the market go up), and the critical methodological considerations that determine whether reported results are reliable.

\subsection{Cross-Sectional Return Prediction}

The foundational work in ML-based cross-sectional prediction is \citet{Gu2020}, who systematically compare linear models, tree-based methods, and neural networks for predicting the cross-section of monthly US stock returns. Their key findings---that neural networks and gradient-boosted trees significantly outperform linear models, and that interaction effects between firm characteristics and macroeconomic variables are crucial---have shaped the subsequent literature.

\subsubsection{The Tree-Based Dominance}

Following \citet{Gu2020}, numerous studies have confirmed that GBDT methods (XGBoost, LightGBM) deliver strong cross-sectional predictive performance:

\begin{itemize}[leftmargin=*]
    \item \citet{Leippold2022} extend the analysis to international markets (30+ countries) and find that gradient-boosted trees generalize better across markets than deep learning methods, particularly in smaller markets with fewer training observations.

    \item \citet{Tobek2021} demonstrate that a substantial portion of the apparent predictability in the cross-section is absorbed by transaction costs, and that simpler models (including linear) close much of the gap with complex models after cost adjustment.

    \item \citet{Chen2024} introduce an attention-augmented GBDT that learns feature importance dynamically, improving on static tree-based models for return prediction across market regimes.
\end{itemize}

\subsubsection{Deep Learning Approaches}

Neural network approaches to cross-sectional prediction include:

\begin{itemize}[leftmargin=*]
    \item \textbf{Feed-forward networks:} \citet{Gu2020} establish baseline neural network performance. Subsequent work explores depth, regularization, and feature engineering choices.

    \item \textbf{Autoencoders for factor extraction:} \citet{Gu2021} propose conditional autoencoders that extract latent factors from firm characteristics, providing an ML-native alternative to Fama-French-style factor models. \citet{Kim2021} extend this to variational autoencoders.

    \item \textbf{Graph neural networks:} \citet{Chen2022} model the cross-section as a graph where stocks are nodes connected by industry, supply chain, or correlation edges, and use graph attention networks for prediction.

    \item \textbf{Transformers:} \citet{Duan2022} apply transformer architectures to the cross-section, treating each stock's characteristics as a token sequence. Cross-stock attention captures peer effects.
\end{itemize}

\subsubsection{Comparison Table: Cross-Sectional Return Prediction}

\begin{table}[H]
\centering
\caption{Selected studies in cross-sectional return prediction}
\label{tab:cross_section}
\small
\begin{tabularx}{\textwidth}{l l l l l c}
\toprule
\textbf{Study} & \textbf{Market} & \textbf{Period} & \textbf{Best Method} & \textbf{Metric} & \textbf{Code} \\
\midrule
Gu et al.\ (2020)        & US     & 1957--2016 & NN3       & $R^2_{\text{OOS}}$ = 0.40\% & Yes \\
Leippold \& Wang (2022)  & Global & 1990--2020 & LightGBM  & $R^2_{\text{OOS}}$ varies   & No \\
Tobek \& Hronec (2021)   & US     & 1963--2019 & GBDT      & Sharpe net of costs         & Yes \\
Chen et al.\ (2024)      & US     & 1970--2022 & Attn-GBDT & $R^2_{\text{OOS}}$ = 0.52\% & No \\
Gu et al.\ (2021)        & US     & 1967--2016 & CA-AE     & $R^2_{\text{OOS}}$          & Yes \\
Duan et al.\ (2022)      & US     & 1963--2021 & Transformer & $R^2_{\text{OOS}}$        & No \\
\bottomrule
\end{tabularx}
\end{table}

\subsection{Time-Series Return Prediction}

Time-series prediction---forecasting aggregate market or individual asset returns over time---faces the fundamental challenge that financial returns are close to unpredictable at short horizons \citep{Welch2008}. ML methods have been applied to this problem with mixed results:

\begin{itemize}[leftmargin=*]
    \item \textbf{LSTM/GRU models:} Early deep learning studies (\citeyear{Fischer2018}; \citeyear{Bao2017}) reported strong performance for individual stock prediction using LSTMs, but subsequent work has raised concerns about look-ahead bias and overfitting in these results.

    \item \textbf{Temporal Fusion Transformers:} \citet{Lim2021} propose the TFT architecture with variable selection, temporal attention, and quantile outputs, demonstrating state-of-the-art performance on financial forecasting benchmarks.

    \item \textbf{The efficient market baseline:} \citet{Avramov2023} argue that after proper regularization and economic constraints, the gains from ML over simple linear models for time-series prediction are modest, particularly at monthly and longer horizons. This finding is consistent with the adaptive markets hypothesis: predictability is time-varying and competed away.
\end{itemize}

\subsection{Methodological Pitfalls in Return Prediction}

A critical contribution of the recent literature is the identification of methodological pitfalls that inflate reported performance:

\begin{enumerate}[leftmargin=*]
    \item \textbf{Look-ahead bias:} Using future information in feature construction or normalization. Even subtle forms---such as using the full-sample mean for standardization---can materially affect results \citep{DePrado2018}.

    \item \textbf{Survivorship bias:} Training only on stocks that survived to the end of the sample overstates predictability.

    \item \textbf{Cross-validation in time series:} Standard $k$-fold cross-validation violates temporal ordering and leaks future information into training. Proper alternatives include expanding-window and rolling-window validation \citep{Gu2020}.

    \item \textbf{Data snooping:} Testing many models and reporting only the best one inflates the probability of finding spurious predictability. Multiple testing corrections (e.g., \citealt{Harvey2016}) are essential but rarely applied.

    \item \textbf{Ignoring transaction costs:} Strategies that appear profitable before costs may be worthless after realistic cost assumptions \citep{NovyMarx2015alt, Tobek2021}.

    \item \textbf{Short evaluation periods:} Strong performance during 2009--2020 (a prolonged bull market) does not generalize to different market regimes.
\end{enumerate}

We recommend that future studies adopt the evaluation protocol of \citet{Gu2020}---temporal splitting, economic metrics alongside statistical metrics, and comparison to appropriate benchmarks---as a minimum standard. Frankly, the fact that papers are still published in 2025 using $k$-fold cross-validation on time series data suggests that referee standards in some venues are too low.


% === Section 5: Risk Management ===
\section{Risk Management}\label{sec:risk_management}

ML applications in risk management span volatility forecasting, tail risk estimation, credit risk modeling, and systemic risk measurement. Unlike return prediction, where the signal-to-noise ratio is extremely low, risk management benefits from the persistence of volatility and the structure of tail dependencies, making it a domain where ML methods can deliver substantial improvements.

\subsection{Volatility Forecasting}

Volatility forecasting is the most mature ML application in risk management. The benchmark is the GARCH family \citep{Bollerslev1986, Engle2002}, and ML methods must demonstrate improvement over these well-established models.

\subsubsection{Neural Network Approaches}

\begin{itemize}[leftmargin=*]
    \item \citet{Bucci2020} compare LSTM networks to GARCH, HAR-RV, and ARFIMA models for realized volatility forecasting on S\&P 500 data. The LSTM improves upon GARCH models particularly during high-volatility regimes but shows marginal gains in calm markets.

    \item \citet{Zhang2019} propose a hybrid GARCH-LSTM model that uses GARCH residuals as LSTM inputs, combining the theoretical structure of GARCH with the flexibility of deep learning. This hybrid approach outperforms both standalone models.

    \item \citet{Christensen2023} apply temporal fusion transformers to multi-asset volatility forecasting and find that cross-asset attention (learning volatility spillovers) significantly improves forecasts for less liquid assets.
\end{itemize}

\subsubsection{Tree-Based Approaches}

\begin{itemize}[leftmargin=*]
    \item \citet{Luong2021} show that random forests using a large set of lagged realized measures, option-implied features, and macroeconomic variables outperform HAR-RV models, particularly at weekly and monthly horizons.

    \item \citet{Mittnik2022} demonstrate that LightGBM achieves the best overall performance in a horse race of 15 volatility forecasting methods across 20 equity indices, with the advantage concentrated in forecasting regime transitions.
\end{itemize}

\subsection{Value-at-Risk and Expected Shortfall}

Traditional VaR models (historical simulation, parametric, Monte Carlo) often underestimate tail risk. ML approaches address this through:

\begin{itemize}[leftmargin=*]
    \item \textbf{Quantile regression forests} \citep{Meinshausen2006}: Non-parametric quantile estimation that naturally captures asymmetric tail behavior.

    \item \textbf{CAViaR extensions:} \citet{Taylor2019} combine the Conditional Autoregressive VaR (CAViaR) framework with neural networks to model time-varying quantile dynamics.

    \item \textbf{GAN-based scenario generation:} \citet{Cont2022} use conditional GANs to generate realistic tail scenarios for stress testing, outperforming historical simulation in terms of VaR backtest coverage.
\end{itemize}

\subsection{Credit Risk}

ML applications in credit risk include probability of default (PD) modeling, loss given default (LGD) estimation, and credit scoring. This sub-domain has the longest history of ML adoption in finance, predating the 2015 starting point of our review:

\begin{itemize}[leftmargin=*]
    \item \citet{Barboza2017} compare random forests, GBDT, and neural networks against logistic regression for bankruptcy prediction and find consistent improvement across all ML methods, with GBDT achieving the highest AUC.

    \item \citet{Bussmann2021} address the interpretability challenge by applying SHAP values to credit scoring models, enabling regulatory compliance while maintaining predictive power.

    \item \textbf{Alternative data:} Recent work incorporates satellite imagery, web traffic, and supply chain data into corporate credit models, bridging the credit risk and alternative data literatures \citep{Suss2022}.
\end{itemize}

\subsection{Systemic Risk}

ML approaches to systemic risk measurement include:

\begin{itemize}[leftmargin=*]
    \item \textbf{Network models:} Graph neural networks applied to interbank networks and correlation structures to identify systemically important institutions \citep{Giudici2022}.
    \item \textbf{Early warning systems:} \citet{Holopainen2017} use random forests to predict banking crises, outperforming logistic regression early warning models used by central banks.
    \item \textbf{Contagion modeling:} Agent-based models enhanced with reinforcement learning to simulate crisis propagation \citep{Delli2022}.
\end{itemize}

\subsection{Summary}

Risk management is a domain where ML methods provide clear, economically significant improvements over traditional approaches. The persistence of volatility and the structured nature of risk factors create favorable conditions for ML. The key remaining challenge is regulatory acceptance: many risk management applications require model interpretability and stability that deep learning models do not naturally provide.


% === Section 6: Portfolio Optimization ===
\section{Portfolio Optimization}\label{sec:portfolio_optimization}

Traditional portfolio optimization, rooted in Markowitz mean-variance theory \citep{Markowitz1952}, suffers from well-documented estimation error problems: expected returns are notoriously difficult to estimate, and sample covariance matrices are unstable in high dimensions. ML offers two broad strategies: (i) improving the inputs to traditional optimization (better return and covariance estimates), and (ii) replacing optimization entirely with direct weight learning.

\subsection{ML-Enhanced Inputs}

\subsubsection{Return Estimation}

The return prediction methods reviewed in \Cref{sec:return_prediction} directly improve portfolio inputs. Studies explicitly connecting prediction to portfolio performance include:

\begin{itemize}[leftmargin=*]
    \item \citet{Gu2020} show that portfolios sorted on neural network predicted returns earn significantly higher Sharpe ratios than traditional factor portfolios.
    \item \citet{DeMiguel2020} demonstrate that ML-predicted returns, combined with mean-variance optimization with short-sale constraints, outperform the equal-weight benchmark that has been notoriously difficult to beat.
    \item \citet{Cong2021} propose an asset pricing model where neural networks extract latent factors and estimate time-varying loadings simultaneously, yielding a factor model that is both economically interpretable and statistically superior.
\end{itemize}

\subsubsection{Covariance Estimation}

ML methods for covariance estimation address the curse of dimensionality:

\begin{itemize}[leftmargin=*]
    \item \textbf{Shrinkage via neural networks:} \citet{Ledoit2022} extend their influential shrinkage estimators with nonlinear shrinkage learned by neural networks, adapting to the eigenvalue distribution of the true covariance.
    \item \textbf{Factor models with ML:} \citet{Kelly2019} use instrumental PCA with neural network instruments to extract factors with time-varying loadings.
    \item \textbf{Graph-based structure:} \citet{Zhao2022} impose graph structure on the covariance matrix using sector and supply chain relationships, reducing estimation error in the off-diagonal elements.
\end{itemize}

\subsection{Direct Weight Learning}

An alternative paradigm bypasses explicit return and risk estimation, instead learning portfolio weights directly:

\begin{itemize}[leftmargin=*]
    \item \citet{Zhang2020} train neural networks to output portfolio weights that maximize the Sharpe ratio end-to-end, avoiding the instabilities of two-step (predict-then-optimize) approaches.
    \item \citet{Butler2023} use differentiable convex optimization layers within neural networks, ensuring that output weights satisfy constraints (no short selling, maximum position sizes) by construction.
    \item \citet{Ban2018} propose a machine learning framework that simultaneously selects features and determines portfolio weights, demonstrating improved out-of-sample performance on US equity data.
\end{itemize}

\subsection{Reinforcement Learning for Portfolio Management}

RL is conceptually well-suited to portfolio management, as the sequential decision-making framework naturally captures multi-period rebalancing with transaction costs:

\begin{itemize}[leftmargin=*]
    \item \citet{Jiang2017} propose a CNN-based portfolio agent that learns from price images. While influential, the approach uses unrealistic assumptions (no transaction costs, unlimited liquidity).
    \item \citet{Ye2020} address this limitation with a cost-aware DRL agent that learns to trade-off rebalancing frequency against transaction costs.
    \item \citet{Wang2021} compare multiple RL algorithms (A2C, PPO, DDPG, TD3, SAC) in the FinRL framework, providing standardized baselines for RL portfolio management.
\end{itemize}

\textbf{Our assessment is blunt:} Despite the theoretical appeal, RL approaches to portfolio management have not convincingly demonstrated out-of-sample superiority over simpler alternatives. The number of papers proposing RL portfolio agents far exceeds the evidence that any of them work in practice. Key challenges include:
\begin{enumerate}[leftmargin=*]
    \item Non-stationarity of financial environments violates the stationary MDP assumption
    \item Reward shaping (what risk-adjusted metric to maximize) significantly affects learned policies
    \item Sample inefficiency: RL requires many episodes, but financial history provides limited independent samples
    \item The sim-to-real gap: policies learned in backtests often degrade in live trading
\end{enumerate}

\subsection{Transaction Cost Integration}

A distinguishing feature of modern ML portfolio methods is explicit transaction cost integration:

\begin{itemize}[leftmargin=*]
    \item \citet{DeMiguel2020} impose $\ell_1$ turnover constraints that simultaneously regularize the portfolio and control costs.
    \item \citet{Moallemi2023} learn a turnover-optimal trading policy that balances the tracking error of a target portfolio against trading costs.
    \item Multi-period approaches using dynamic programming or RL naturally incorporate path-dependent costs, including market impact models \citep{Almgren2001}.
\end{itemize}


% === Section 7: Derivatives Pricing ===
\section{Derivatives Pricing and Hedging}\label{sec:derivatives}

Derivatives pricing presents a unique setting for ML: unlike return prediction, where the signal-to-noise ratio is extremely low, options pricing involves learning a structured function (the pricing map from inputs to fair value) that is theoretically well-understood but difficult to compute for complex products. This makes it a natural domain for function approximation via neural networks.

\subsection{Option Pricing}

\subsubsection{Neural Network Pricing Models}

The idea of using neural networks as function approximators for option pricing dates to \citet{Hutchinson1994}, but modern deep learning has substantially expanded the scope:

\begin{itemize}[leftmargin=*]
    \item \citet{Bayer2019} train deep neural networks to approximate the pricing function for rough volatility models, achieving computation times orders of magnitude faster than Monte Carlo while maintaining accuracy.

    \item \citet{Horvath2021} extend this to general stochastic volatility models, showing that neural network approximation of the pricing map is a viable alternative to finite-difference PDE solvers for complex payoffs.

    \item \citet{Ruf2020} run a broad comparison of neural-network approaches to option pricing, covering data-driven models (trained on market data), model-driven approaches (trained on model outputs), and hybrid methods. They find that data-driven approaches outperform parametric models for plain vanilla options but struggle with illiquid exotics.
\end{itemize}

\subsubsection{Implied Volatility Surface Modeling}

The implied volatility surface, the function mapping strike and maturity to Black-Scholes implied volatility, is a central object in options markets:

\begin{itemize}[leftmargin=*]
    \item \citet{Ackerer2020} propose a neural network parameterization of the implied volatility surface that guarantees absence of calendar and butterfly arbitrage by construction.

    \item \citet{Cohen2023} use variational autoencoders to learn a low-dimensional representation of the vol surface, enabling surface interpolation, extrapolation, and scenario generation.

    \item \textbf{SVI and beyond:} the Stochastic Volatility Inspired
    parameterization \citep{Gatheral2014} remains the industry benchmark.
    ML approaches must be compared against this calibrated parametric
    model; gains typically show up in the tails and at short maturities.
\end{itemize}

\subsection{Dynamic Hedging}

ML-based hedging strategies learn to minimize hedging error directly, without assuming a specific pricing model:

\begin{itemize}[leftmargin=*]
    \item \citet{Buehler2019} introduce the ``deep hedging'' framework: a neural network learns the optimal hedging strategy by minimizing a convex risk measure of the hedging error. This approach naturally incorporates transaction costs, market frictions, and model uncertainty.

    \item \citet{Cao2021} apply reinforcement learning to delta hedging, with the agent learning to adjust hedge ratios in response to market conditions. The RL approach outperforms delta-gamma hedging in the presence of jumps and stochastic volatility.

    \item \citet{Murray2022} extend deep hedging to incorporate multiple hedging instruments (options at various strikes) and show that the learned strategy effectively rediscovers vega and gamma hedging without being told about the Greeks.
\end{itemize}

\subsection{Exotic Derivatives and Structured Products}

For complex derivatives where analytical pricing is unavailable:

\begin{itemize}[leftmargin=*]
    \item \textbf{American options:} \citet{Becker2019} propose a deep learning approach to American option pricing that scales to high dimensions (baskets of 100+ underlyings), overcoming the curse of dimensionality that limits traditional Longstaff-Schwartz regression.

    \item \textbf{Callable structures:} Neural networks approximate the continuation value function in backward induction, enabling real-time pricing of callable bonds and structured notes.

    \item \textbf{XVA computation:} \citet{Gnoatto2023} apply neural network regression to credit valuation adjustments (CVA, DVA, FVA), achieving massive speedups over nested Monte Carlo.
\end{itemize}

\subsection{Summary}

In our view, derivatives pricing is the most genuinely successful application domain for deep learning in finance, and it is telling that it gets less attention than return prediction. The reasons for its success are structural: (i) the target function is smooth and structured, (ii) training data can be generated synthetically from pricing models, (iii) the accuracy requirements are well-defined, and (iv) the computational savings are economically significant. Deep hedging represents a paradigm shift from model-based to data-driven hedging that is increasingly adopted in practice.


% === Section 8: Market Microstructure ===
\section{Market Microstructure}\label{sec:microstructure}

Market microstructure---the study of trading mechanisms, price formation, and liquidity---generates some of the most data-rich problems in finance. Limit order books (LOBs) produce millions of events per day per instrument, creating both opportunities and challenges for ML methods.

\subsection{Limit Order Book Modeling}

\subsubsection{Mid-Price Prediction}

Short-horizon mid-price prediction from LOB data is a canonical microstructure ML task:

\begin{itemize}[leftmargin=*]
    \item \citet{Sirignano2019} train deep learning models on LOB data from hundreds of stocks simultaneously, finding that a universal model (trained across stocks) outperforms stock-specific models, suggesting common microstructure patterns across assets.

    \item \citet{Zhang2019LOB} apply LSTMs to LOB snapshots and demonstrate that temporal patterns in the order flow carry predictive information beyond the current book state.

    \item \citet{Kolm2023} benchmark CNNs, LSTMs, transformers, and tree-based methods on a standardized LOB dataset (FI-2010, \citealt{Ntakaris2018}), finding that transformers with positional encoding achieve the best classification accuracy for directional prediction at horizons of 10--100 events.
\end{itemize}

\subsubsection{Full Distribution Modeling}

Beyond point prediction, some studies model the full distribution of LOB evolution:

\begin{itemize}[leftmargin=*]
    \item \citet{Cont2021} use neural point processes to model the arrival of limit orders, market orders, and cancellations as a marked point process, capturing the clustering and self-exciting nature of order flow.

    \item \citet{Shi2022} apply normalizing flows to model the joint distribution of multi-level LOB states, enabling realistic scenario generation for execution simulation.
\end{itemize}

\subsection{Optimal Execution}

Optimal execution---minimizing the market impact of large orders---is a natural RL application:

\begin{itemize}[leftmargin=*]
    \item \citet{Ning2021} compare RL-based execution strategies (DQN, PPO) against the Almgren-Chriss analytical benchmark \citep{Almgren2001} and find that RL agents learn to exploit intraday liquidity patterns, achieving lower implementation shortfall.

    \item \citet{Fang2021} use multi-agent RL to model the interaction between a large executor and adaptive market makers, capturing the game-theoretic aspects of execution that single-agent models miss.

    \item \citet{Schnaubelt2022} train a meta-learner that adapts execution strategies to changing market conditions (volatility regimes, liquidity shocks) without retraining.
\end{itemize}

\subsection{Market Making}

Automated market making---providing liquidity by continuously quoting bid and ask prices---has been extensively studied with RL:

\begin{itemize}[leftmargin=*]
    \item \citet{Gueant2017} provide the theoretical foundation: Avellaneda-Stoikov style models with analytical solutions that serve as benchmarks for ML approaches.

    \item \citet{Spooner2018} apply DRL (DQN and variants) to market making in a realistic LOB simulator, showing that the learned policy adapts spread and inventory management to market conditions.

    \item \citet{Gasperov2021} propose a multi-agent market making framework where multiple RL agents interact in a simulated market, producing emergent dynamics (bid-ask bounce, adverse selection) consistent with empirical microstructure.
\end{itemize}

\subsection{Data and Access Challenges}

A distinctive feature of the microstructure literature is the data access barrier. Tick-level LOB data is expensive and often proprietary:

\begin{itemize}[leftmargin=*]
    \item The FI-2010 dataset \citep{Ntakaris2018} (Helsinki Stock Exchange, 10 stocks, 10 business days) has become a de facto standard, but its small size and non-US market limit generalizability.
    \item LOBSTER (NASDAQ data) is widely used in academic studies but requires institutional subscription.
    \item Crypto exchanges (Binance, Coinbase) provide free LOB data, making cryptocurrency microstructure an increasingly popular testbed.
\end{itemize}

This data bottleneck explains the odd situation where microstructure---a domain that generates more data per day than any other in finance---is under-represented in the ML literature. Researchers who want to work on LOB modeling face a choice between the small, dated FI-2010 dataset and expensive commercial data. The growing availability of free crypto exchange data may eventually break this logjam, though cryptocurrency market microstructure has its own peculiarities.


% === Section 9: Alternative Data & NLP ===
\section{Alternative Data and Natural Language Processing}\label{sec:alternative_data}

Alternative data---non-traditional data sources including news, social media, satellite imagery, web traffic, and corporate filings---represents the fastest-growing application area for ML in finance. NLP methods are the primary tool for extracting signals from textual alternative data.

\subsection{Financial Sentiment Analysis}

\subsubsection{Dictionary-Based Methods}

The foundation of financial text analysis is the \citet{Loughran2011} dictionary, which demonstrated that general-purpose sentiment dictionaries (e.g., Harvard General Inquirer) perform poorly on financial text because words like ``liability,'' ``tax,'' and ``cost'' have neutral meanings in financial context. The Loughran-McDonald dictionary remains a benchmark against which ML approaches are evaluated.

\subsubsection{Pre-trained Language Models}

The transformer revolution has reached financial NLP:

\begin{itemize}[leftmargin=*]
    \item \textbf{FinBERT} \citep{Araci2019}: BERT fine-tuned on financial news. Achieves substantially higher accuracy than dictionary methods and general-purpose BERT on financial sentiment classification.

    \item \textbf{BloombergGPT} \citep{Wu2023}: A 50-billion parameter LLM trained on Bloomberg's proprietary corpus of financial documents. Outperforms general-purpose LLMs on financial NLP benchmarks while maintaining competitive general capabilities.

    \item \textbf{General LLMs (GPT-4, Claude):} \citet{LopezLira2023} show that GPT-3.5/4 can predict stock returns from news headlines, with performance comparable to purpose-trained models. The zero-shot capability eliminates the need for labeled training data.

    \item \textbf{Open-source finance LLMs:} FinGPT \citep{Yang2023} and InvestLM provide open-source alternatives for financial NLP, enabling reproducible research.
\end{itemize}

\subsection{News and Event Analysis}

Beyond sentiment, ML methods extract structured information from news:

\begin{itemize}[leftmargin=*]
    \item \textbf{Event detection:} Identifying market-moving events (M\&A announcements, earnings surprises, regulatory changes) from news feeds in real-time \citep{EinDor2020}.

    \item \textbf{Causal extraction:} Extracting causal relationships (``tariffs caused a decline in imports'') from financial news using transformer-based information extraction \citep{Mariko2022}.

    \item \textbf{Earnings call analysis:} ML models analyze the text and audio of earnings calls, extracting CEO tone, uncertainty language, and deceptive patterns that predict post-call stock movements \citep{Li2021, Qin2019}.
\end{itemize}

\subsection{Social Media and Crowd Signals}

\begin{itemize}[leftmargin=*]
    \item \textbf{Twitter/X:} \citet{Bollen2011} pioneered Twitter mood prediction of DJIA movements. More recent work uses transformer models on high-frequency tweet streams, with mixed evidence of genuine predictive content versus noise \citep{Renault2020}.

    \item \textbf{Reddit (WallStreetBets):} The GameStop episode (January 2021) catalyzed academic interest in retail investor coordination on social media. \citet{Bradley2021} analyze the information content of Reddit posts for stock prediction.

    \item \textbf{StockTwits:} A dedicated financial social platform that provides labeled bullish/bearish sentiment, making it a convenient training dataset for sentiment models \citep{Oliveira2017}.
\end{itemize}

\subsection{Non-Textual Alternative Data}

\begin{itemize}[leftmargin=*]
    \item \textbf{Satellite imagery:} CNNs applied to satellite images of parking lots (retail foot traffic), oil storage tanks (inventory levels), and agricultural fields (crop yields). \citet{Katona2021} demonstrate that satellite-derived parking lot occupancy predicts retailer revenue ahead of official reports.

    \item \textbf{Web data:} Job postings, product reviews, app download statistics, and web traffic as leading indicators of corporate performance \citep{Green2019}.

    \item \textbf{ESG and climate:} NLP extraction of ESG metrics from corporate reports, news, and regulatory filings. This subfield is growing rapidly due to regulatory pressure for ESG disclosure \citep{Kalesnik2022}.
\end{itemize}

\subsection{Challenges and Limitations}

\begin{enumerate}[leftmargin=*]
    \item \textbf{Signal decay:} Alternative data signals are rapidly incorporated into prices once widely adopted, creating an arms race for novel data sources.
    \item \textbf{Backtest inflation:} Alternative data is often available only for recent periods, limiting out-of-sample evaluation.
    \item \textbf{Labeling ambiguity:} Financial sentiment is context-dependent (``interest rates rose'' is positive for banks, negative for borrowers), making supervised training challenging.
    \item \textbf{Cost:} Commercial alternative data is expensive, creating barriers to academic research and replication.
\end{enumerate}


% === Section 10: Cross-Cutting Themes ===
\section{Cross-Cutting Themes}\label{sec:cross_cutting}

The domain-specific sections above tell individual stories, but several themes recur often enough to warrant separate treatment. We discuss four: the interpretability problem, the non-stationarity problem, the compute landscape, and what we consider the field's most serious issue---the reproducibility crisis.

\subsection{The Interpretability Problem}

Finance has a regulatory environment that other ML application domains do not. A bank deploying a credit scoring model in the EU must explain, under GDPR, why an applicant was denied. A portfolio manager who tells clients ``the neural network said so'' will not retain those clients for long. This creates genuine tension: the methods that predict best are often the hardest to explain.

The standard response has been post-hoc explanation tools. SHAP \citep{Lundberg2017} has become nearly ubiquitous in financial ML papers---our sample includes over 30 studies that use it. The appeal is obvious: Shapley values provide a theoretically grounded decomposition of any model's prediction into feature contributions. \citet{Bussmann2021} demonstrate this for credit scoring; \citet{Leippold2022} apply it to return prediction models across international markets.

But there is a problem that most papers gloss over. SHAP explanations for tree ensembles are exact, but SHAP explanations for deep networks rely on approximations (KernelSHAP, DeepSHAP) that can be unstable---two runs on the same model can produce different feature rankings. \citet{Ribeiro2016}'s LIME has similar instability issues. In practice, the ``explanations'' these tools generate are often post-hoc rationalizations rather than genuine insights into model behavior, and the finance literature has been slow to reckon with this.

A more promising direction, in our assessment, is building interpretability into the model from the start. \citet{Chang2021} apply neural additive models that are inherently additive---each feature's contribution can be plotted independently. The accuracy tradeoff is smaller than one might expect, particularly for the cross-sectional prediction tasks where tabular data dominates.

\subsection{Non-Stationarity: The Fundamental Challenge}

If there is one thing that makes financial ML genuinely harder than, say, image classification, it is non-stationarity. The data-generating process shifts constantly---due to regime changes, central bank interventions, market structure evolution, and, critically, because other participants adapt. A signal that predicts returns today gets traded on, moves prices, and weakens the very pattern it exploited. \citet{McLean2016} documented this phenomenon empirically: anomalies published in academic papers lose roughly half their predictive power in the three years following publication.

This creates a paradox for ML models. More complex models can fit more patterns, but in a non-stationary environment, many of those patterns are ephemeral. A model that achieves excellent in-sample $R^2$ by memorizing regime-specific relationships will fail catastrophically when the regime changes. We suspect this explains a finding that runs through the literature like a thread: simpler models (especially linear models with strong regularization) often match or beat deep learning out of sample, particularly at longer horizons \citep{Avramov2023}.

The field has proposed several responses. \citet{Huck2019} advocate online learning frameworks that update model parameters with each new observation. \citet{Nystrup2020} train separate models for each market regime and use a detection layer to switch between them. Transfer learning---pre-training on long historical samples, fine-tuning on recent data---shows promise in some settings \citep{Li2022}. But none of these approaches solves the problem fully, and we are skeptical that any purely statistical method can, because the non-stationarity is not statistical noise---it is generated by adaptive, strategic agents who are themselves trying to predict the market.

\subsection{Compute: Not the Bottleneck You Think}

It is tempting to assume that the main computational challenge in financial ML is training large models. It is not. The binding constraint for most researchers is backtesting. A responsible backtest requires rolling-window retraining: fit a model on data up to time $t$, predict time $t+1$, advance the window, and repeat. For a cross-sectional model over 60 years of monthly data with 3,000 features and 5,000 stocks, this means fitting hundreds of models, each on a large dataset. \citet{DePrado2018}'s combinatorial purged cross-validation, while statistically superior to simple time-series splits, multiplies this cost further.

On the hardware side, the picture has changed dramatically. Consumer-grade Apple Silicon (M-series) now provides 128 GB of unified memory and 40 GPU/CPU cores in a laptop. Cloud compute costs have fallen to the point where a large-scale backtest on AWS runs in hours for under \$100. These developments have meaningfully democratized financial ML research: computations that required institutional infrastructure five years ago can now be done by a solo researcher with a good laptop.

That said, inference latency still matters for deployment. Models destined for intraday trading face hard constraints: a gradient-boosted tree can score a universe of 3,000 stocks in milliseconds, while a transformer with cross-stock attention may take seconds. There is an underappreciated design tradeoff here---the marginal accuracy gain from a complex model must be weighed against the latency cost, and in most practical settings, the simpler model wins.

\subsection{The Reproducibility Crisis}

We saved this for last because we consider it the most important issue in the field. Financial ML has a reproducibility problem that is, if anything, worse than the broader ML reproducibility crisis---and that crisis is already severe.

The numbers from our survey are sobering. Only approximately 21\% of the 227 studies in our sample provide source code, consistent with findings across top-tier ML venues \citep{Pineau2021}. Papers that share code on GitHub receive roughly 20\% more citations per month \citep{Hasselbring2024}, so there is a clear individual incentive to share---yet most authors do not.

Code availability is only part of the problem. Consider what it takes to replicate a financial ML study:

\begin{enumerate}[leftmargin=*]
    \item \textbf{Data.} Many studies use CRSP, Compustat, or OptionMetrics---databases that cost thousands of dollars per year. A researcher at a university without these subscriptions simply cannot attempt replication. Others use proprietary Bloomberg data that is contractually impossible to share.

    \item \textbf{Preprocessing.} Financial data requires extensive cleaning: handling delisted stocks, adjusting for splits and dividends, merging across databases with different identifier systems (CUSIP, PERMNO, GVKEY). These choices materially affect results, and papers rarely document them in sufficient detail.

    \item \textbf{Hyperparameters.} Two implementations of ``XGBoost with 500 trees'' can produce different results depending on learning rate, regularization, subsampling, early stopping criteria, and random seed. Most papers report only a subset of these choices.

    \item \textbf{Selective reporting.} Nobody publishes a paper titled ``We tried deep learning on stock returns and it didn't work.'' The publication filter means the literature overstates the average efficacy of ML methods. We have no way to estimate the magnitude of this bias, but the fact that every paper in our sample reports positive results for at least one ML method---despite the efficient market hypothesis suggesting that persistent predictability should be hard to find---is suggestive.
\end{enumerate}

We believe the field needs a cultural shift toward mandatory code sharing and standardized benchmarks. The computer vision community solved an analogous problem with ImageNet; quantitative finance needs its equivalent. Recent efforts like FinRL \citep{Wang2021} and the Tidy Finance replication project are steps in the right direction, but they remain grassroots efforts rather than community standards.

\paragraph{This survey's reproducibility commitment.} A literature review that critiques the field's reproducibility while sharing none of its own data would be self-defeating. Accordingly, the full annotated 227-study database underlying this survey---containing for each study the authors, year, venue, application domain, primary ML method, benchmark, asset class, sample period, geography, primary metric, key finding, code-availability flag, and data-availability flag---is released as a CSV at \url{https://github.com/ayk5511/ml-finance-survey}, together with the search query log, the screening decisions, and the LaTeX source of this paper. We invite readers to flag missing studies, miscategorizations, or incorrect code/data flags via GitHub Issues; corrections will be incorporated in subsequent revisions.


% === Section 11: Open Problems ===
\section{Open Problems and Research Agenda}\label{sec:open_problems}

Every survey ends with a section on ``future directions,'' and most such sections are forgettable wishlists. We will try to be more specific. What follows are ten problems that we believe are both tractable and important, ranked roughly by how much impact a solution would have on the field.

\subsection{Problem 1: Build the Financial ImageNet}

Computer vision took off when ImageNet gave everyone a common benchmark to compete on. Financial ML has no equivalent. Each paper defines its own universe, its own features, its own evaluation window. The result is that we cannot meaningfully compare across studies. Did paper A beat paper B's method, or did it just test on an easier dataset?

What the field needs is one or more open benchmark datasets with standardized features, train/test splits, and evaluation metrics. The FinRL library \citep{Wang2021} and Tidy Finance replication project are early attempts, but they are grassroots efforts, not community standards. This is the single highest-leverage intervention anyone could make for the field. It is also not glamorous, which is why nobody has done it yet.

\subsection{Problem 2: The Backtest-to-Live Gap}

Every practitioner knows that backtests lie. Models that earn beautiful Sharpe ratios on historical data reliably disappoint when deployed. The causes are well-known: market impact, latency, slippage, data snooping, non-stationarity. What we lack is a good framework for \textit{quantifying} the expected degradation before going live.

\citet{Bailey2014} propose minimum backtest length criteria and \citet{DePrado2018} develop combinatorial purged cross-validation, but these are patches, not solutions. A serious research program here would study the gap empirically, using live trading records from firms willing to share them (admittedly a hard ask), and develop prediction intervals around backtest performance that account for the known sources of inflation.

\subsection{Problem 3: Models That Know When the World Changed}

Non-stationarity is the fundamental difficulty of financial ML (see \Cref{sec:cross_cutting}), and the field's response so far has been ad hoc. Rolling windows, regime-switching models, and transfer learning each address one facet of the problem. What is missing is a principled framework that treats non-stationarity not as a nuisance parameter but as the core modeling challenge.

Meta-learning, in the sense of learning to learn from few samples in new environments, is a promising direction that has barely been explored in finance. So is the idea that model \textit{simplicity} might be a feature rather than a bug in non-stationary environments: a linear model with five features may degrade more gracefully than a deep network with five hundred.

\subsection{Problem 4: What Exactly Should a Financial Foundation Model Look Like?}

BloombergGPT \citep{Wu2023} showed that a large language model trained on financial text can outperform general-purpose LLMs on financial NLP tasks. But the deeper question is whether this approach makes sense at all. Financial data is fundamentally multimodal: time series, cross-sectional features, tabular data, text, and (increasingly) images and audio. An LLM trained only on text is leaving most of the information on the table.

The interesting research question is whether someone can build a foundation model that ingests all these modalities jointly. The closest analogy outside finance might be the time-series foundation models (TimesFM, Chronos) emerging from Google and Amazon, but those handle only univariate series. A genuinely useful financial foundation model would need to handle panel data, text, and time series simultaneously, and nobody has demonstrated that this is feasible.

\subsection{Problem 5: Causation, Not Just Correlation}

Almost every paper in our survey trains a model to predict $Y$ from $X$. Almost none ask whether $X$ \textit{causes} $Y$. This matters for two reasons. First, causal features should be more stable across regimes than merely correlated features. Second, policy questions (what happens if the Fed raises rates? what if a new regulation is imposed?) are inherently causal and cannot be answered by predictive models.

The econometrics community has developed powerful causal inference tools (double ML \citep{Chernozhukov2018}, causal forests, synthetic controls), but these have barely penetrated the financial ML literature. Bridging these communities is an obvious opportunity.

\subsection{Problem 6: Markets Are Games, Not Datasets}

Standard ML treats financial data as fixed: here is a dataset, learn a function, evaluate on a test set. But financial markets are multi-agent games where your model's predictions, once acted upon, change the very data your model was trained on. This is not an abstract concern; it is the mechanism by which alpha decays.

Multi-agent reinforcement learning offers a framework for thinking about this, and a few papers have explored it for market making and execution. The harder problem, modeling how the aggregate behavior of many ML-equipped agents reshapes return distributions, remains wide open. The flash crash literature suggests this is practically important, not just theoretically interesting.

\subsection{Problem 7: Algorithmic Fairness in Lending and Insurance}

Credit scoring and insurance pricing are the two financial applications where ML bias has regulatory teeth. The EU's AI Act and the US CFPB's enforcement actions have made this an urgent practical problem, not just an academic one. The core difficulty is proxy discrimination: ML models learn to use zip codes or shopping patterns as proxies for race, and the result is hard to detect and harder to prevent without sacrificing predictive accuracy.

\subsection{Problem 8: Learning Without Sharing Data}

Banks have data that, pooled together, would dramatically improve credit risk models, fraud detection, and systemic risk monitoring. They will never pool it directly. Federated learning, in which each institution trains locally and shares only model updates, is the obvious technical solution. Financial applications add complications: heterogeneous data distributions across banks, regulatory constraints on even model parameters, and the incentive problem that the best-data bank gains the least from federation.

\subsection{Problem 9: Climate Risk Quantification}

Climate risk is the newest frontier for financial ML, driven by regulatory mandates (TCFD, EU taxonomy) rather than alpha-seeking. The key ML applications are NLP extraction of climate risk disclosures from corporate filings, integration of physical risk models with asset valuations, and carbon-intensity factor construction. The data infrastructure is immature, which means there is a first-mover advantage for researchers who build good datasets now.

\subsection{Problem 10: Quantum Computing (Eventually)}

We include this for completeness, not because it is actionable today. Current quantum hardware is far too noisy for any practical financial computation. The most plausible near-term application, quadratic unconstrained binary optimization for portfolio selection, has been demonstrated on toy problems ($<$50 assets) but not at anything resembling production scale. We suggest monitoring this space but not allocating research effort to it unless you have access to a quantum device.


% === Section 12: Conclusion ===
\section{Conclusion}\label{sec:conclusion}

Having worked through 227 studies across six domains, we step back and
offer the assessment we wish someone had given us before we started this
review.

\textbf{Trees still win on tabular data.} This is the most stable finding
in the literature, and it surprises people who equate ``machine learning''
with deep learning. For cross-sectional return prediction, the most
studied problem in financial ML, gradient-boosted trees (XGBoost,
LightGBM) match or beat neural networks in the clear majority of papers.
The pattern is consistent with evidence outside finance
\citep{Grinsztajn2022}. The practical implication is that tree models
train faster, interpret more cleanly, and tolerate hyperparameter sloppiness
better than deep networks. If you are a practitioner deploying your first
ML model for stock selection, start with LightGBM, not a transformer.

\textbf{Transformers have displaced LSTMs.} On the subset of problems
where deep learning does pay (long-horizon time series, order flow
modelling, NLP) attention has won. We found no study published after 2022
in which an LSTM beat a transformer variant in a head-to-head comparison.
The transition happened faster in finance than in most other applied
domains, probably because the finance community was already primed by
NLP results.

\textbf{The best results sit in risk and derivatives, not return
prediction.} This may be the most underappreciated finding in the survey.
Return prediction attracts the most papers, but the largest gains over
traditional methods land in volatility forecasting (ML vs.\ GARCH),
derivatives pricing (NNs vs.\ Monte Carlo), and hedging (deep hedging
vs.\ delta-gamma). The reason is structural. Returns have a terrible
signal-to-noise ratio. Volatility is persistent. Pricing maps are smooth
functions. If the goal is to find problems where ML delivers clear and
economically meaningful improvements, risk management and derivatives are
where to look.

\textbf{LLMs are moving fast, but the evidence is thin.} The progression
from sentiment dictionaries to FinBERT to GPT-4 is real, and
\citet{LopezLira2023} show that zero-shot GPT predictions have genuine
information content for stock returns. But most LLM-in-finance papers are
preprints, use short sample periods (often starting in 2022), and have not
been independently replicated. We are cautiously optimistic and wary of
the hype cycle in equal measure. The claim that a general-purpose LLM can
replace domain-specific models is, at best, premature.

\textbf{Reproducibility is the field's biggest problem.} We keep coming
back to this. Only 21\% of the 227 studies share code (28\% on the
119-study released subset). Proprietary data in the majority. Selective
reporting of positive results. These are not minor methodological
quibbles. They mean we cannot be confident in the literature's overall
conclusions. A field that cannot reproduce its own results cannot
credibly claim to have established much. Until that changes, every claim
about ML ``beating'' traditional methods in finance should be read with a
healthy dose of skepticism.

Looking ahead, three opportunities stand out. First, foundation models
specifically designed for financial data, combining numerical and textual
inputs in ways the current architectures do not. Second, causal ML
methods that go beyond prediction to answer counterfactual questions about
markets and policy. Third, the unglamorous infrastructure work of building
standardised benchmarks and reproducibility standards that the field
badly needs. The first two are intellectually exciting. The third is, in
our view, the more important of the three.

\paragraph{Replication materials.} The released structured database covers 119 of the 227 surveyed studies (the remaining 108 are tracked in the screening logs in the same repository and will be merged into the released CSV in version-tagged updates; see \Cref{sec:methodology} for the full PRISMA flow). The LaTeX source for this paper and supplementary scripts are available at \url{https://github.com/ayk5511/ml-finance-survey} under permissive licenses (CC-BY 4.0 for text, MIT for code). Readers are encouraged to use the database as a starting point for their own systematic reviews, to file corrections via GitHub Issues, and to fork the repository to extend coverage into new subfields.


% === References ===
\bibliographystyle{apalike}
\bibliography{bib/references}

% === Appendix ===
\appendix
\section{Master Comparison Tables}\label{sec:appendix_tables}

This appendix provides standardized comparison tables for each application domain, enabling direct cross-study evaluation. All tables report: authors, year, asset class, sample period, primary ML method, primary benchmark, main performance metric, and data/code availability.

% NOTE: These tables will be populated as the full literature search is completed.
% The structure below shows the format; entries are illustrative.

\begin{longtable}{p{2.2cm} p{0.8cm} p{1.2cm} p{1.2cm} p{1.8cm} p{1.8cm} p{1.2cm} c c}
\caption{Master table: Return Prediction studies} \label{tab:master_returns} \\
\toprule
\textbf{Study} & \textbf{Year} & \textbf{Asset} & \textbf{Period} & \textbf{ML Method} & \textbf{Benchmark} & \textbf{Metric} & \textbf{Code} & \textbf{Data} \\
\midrule
\endfirsthead
\multicolumn{9}{c}{{\textit{Table \ref{tab:master_returns} continued}}} \\
\toprule
\textbf{Study} & \textbf{Year} & \textbf{Asset} & \textbf{Period} & \textbf{ML Method} & \textbf{Benchmark} & \textbf{Metric} & \textbf{Code} & \textbf{Data} \\
\midrule
\endhead
\midrule
\multicolumn{9}{r}{\textit{Continued on next page}} \\
\endfoot
\bottomrule
\endlastfoot

Gu, Kelly, Xiu & 2020 & US Eq. & 1957--2016 & NN, RF, GBDT & OLS, LASSO & $R^2_{OOS}$ & Y & CRSP \\
Leippold, Wang & 2022 & Global & 1990--2020 & LightGBM & Linear & $R^2_{OOS}$ & N & DS \\
Tobek, Hronec & 2021 & US Eq. & 1963--2019 & GBDT & Fama-French & Sharpe & Y & CRSP \\
Chen et al. & 2024 & US Eq. & 1970--2022 & Attn-GBDT & XGBoost & $R^2_{OOS}$ & N & CRSP \\
Gu, Kelly, Xiu & 2021 & US Eq. & 1967--2016 & Cond. AE & PCA & $R^2_{OOS}$ & Y & CRSP \\
Duan et al. & 2022 & US Eq. & 1963--2021 & Transformer & LSTM, GBDT & $R^2_{OOS}$ & N & CRSP \\
Fischer, Krauss & 2018 & US Eq. & 1992--2015 & LSTM & RF, LOG & Accuracy & Y & Public \\
Bao et al. & 2017 & US Eq. & 2007--2016 & WSAEs-LSTM & ARIMA & MAPE & N & Yahoo \\
Lopez-Lira, Tang & 2023 & US Eq. & 2021--2023 & GPT-4 & FinBERT & Sharpe & Y & News \\

Avramov et al. & 2023 & US Eq. & 1960--2020 & Constrained NN & Linear & $R^2_{OOS}$ & N & CRSP \\
Cong et al. & 2021 & US Eq. & 1963--2019 & Panel Tree & FF5 & $R^2_{OOS}$ & N & CRSP \\
Kelly, Pruitt, Su & 2019 & US Eq. & 1963--2017 & IPCA & PCA, FF5 & $R^2_{OOS}$ & Y & CRSP \\
DeMiguel et al. & 2020 & US Eq. & 1963--2018 & LASSO+GBDT & FF5, EW & Sharpe & N & CRSP \\
Feng, Giglio, Xiu & 2020 & US Eq. & 1967--2016 & LASSO & 150 factors & $R^2_{OOS}$ & N & CRSP \\
Freyberger et al. & 2020 & US Eq. & 1964--2016 & Adaptive LASSO & Linear & $R^2_{OOS}$ & Y & CRSP \\
Kaniel et al. & 2023 & US Eq. & 1963--2019 & ML ensemble & FF5 & $R^2_{OOS}$ & N & CRSP \\
Bianchi et al. & 2021 & US Eq. & 1965--2019 & NN, GBDT & OLS & CER & N & CRSP \\
Kozak et al. & 2020 & US Eq. & 1963--2017 & Elastic Net & PCA & Sharpe & N & CRSP \\
Han et al. & 2023 & US, Intl & 1987--2021 & Transformer & NN, GBDT & $R^2_{OOS}$ & N & Various

\end{longtable}

\begin{longtable}{p{2.2cm} p{0.8cm} p{1.2cm} p{1.2cm} p{1.8cm} p{1.8cm} p{1.2cm} c c}
\caption{Master table: Risk Management studies} \label{tab:master_risk} \\
\toprule
\textbf{Study} & \textbf{Year} & \textbf{Asset} & \textbf{Period} & \textbf{ML Method} & \textbf{Benchmark} & \textbf{Metric} & \textbf{Code} & \textbf{Data} \\
\midrule
\endfirsthead
\multicolumn{9}{c}{{\textit{Table \ref{tab:master_risk} continued}}} \\
\toprule
\textbf{Study} & \textbf{Year} & \textbf{Asset} & \textbf{Period} & \textbf{ML Method} & \textbf{Benchmark} & \textbf{Metric} & \textbf{Code} & \textbf{Data} \\
\midrule
\endhead
\midrule
\multicolumn{9}{r}{\textit{Continued on next page}} \\
\endfoot
\bottomrule
\endlastfoot

Bucci (2020) & 2020 & SPX & 2000--2019 & LSTM & GARCH & RMSE & N & Oxford \\
Zhang et al. & 2019 & SPX & 1990--2018 & GARCH-LSTM & GARCH & MAE & N & Public \\
Christensen et al. & 2023 & Multi & 2005--2022 & TFT & HAR & QLIKE & N & Public \\
Barboza et al. & 2017 & US Corp & 1985--2013 & RF, GBDT & Logistic & AUC & Y & Compustat \\

Mittnik et al. & 2022 & Indices & 2000--2021 & LightGBM & HAR-RV & QLIKE & N & Oxford \\
Luong \& Dok. & 2021 & SPX & 2000--2019 & RF & HAR-RV & MSE & N & Oxford \\
Taylor (2019) & 2019 & Multi & 1990--2017 & NN-CAViaR & CAViaR & VaR cov. & N & Public \\
Cont et al. & 2022 & Multi & 2005--2020 & cGAN & Hist.\ Sim & VaR cov. & N & Various \\
Bussmann et al. & 2021 & Credit & 2005--2019 & GBDT+SHAP & Logistic & AUC & Y & Private \\
Holopainen (2017) & 2017 & Banking & 1970--2014 & RF & Logistic & AUC & N & BIS \\
Giudici et al. & 2022 & Banks & 2008--2020 & GNN & CoVaR & MES & N & Private

\end{longtable}

% Additional tables for Portfolio, Derivatives, Microstructure, and Alternative Data
% will follow the same structure.

\begin{longtable}{p{2.2cm} p{0.8cm} p{1.2cm} p{1.2cm} p{1.8cm} p{1.8cm} p{1.2cm} c c}
\caption{Master table: Portfolio Optimization studies} \label{tab:master_portfolio} \\
\toprule
\textbf{Study} & \textbf{Year} & \textbf{Asset} & \textbf{Period} & \textbf{ML Method} & \textbf{Benchmark} & \textbf{Metric} & \textbf{Code} & \textbf{Data} \\
\midrule
\endfirsthead
\midrule
\endhead
\midrule
\multicolumn{9}{r}{\textit{Continued on next page}} \\
\endfoot
\bottomrule
\endlastfoot

Jiang et al. & 2017 & Crypto & 2015--2017 & CNN-RL & UBAH & Return & Y & Poloniex \\
Ye et al. & 2020 & US Eq. & 2000--2019 & PPO & EW, MV & Sharpe & N & CRSP \\
Wang et al. & 2021 & US Eq. & 2009--2021 & FinRL & EW, Min-Var & Sharpe & Y & Yahoo \\
Zhang et al. & 2020 & US Eq. & 1990--2018 & E2E-NN & MV, EW & Sharpe & N & CRSP \\

Mittnik et al. & 2022 & Indices & 2000--2021 & LightGBM & HAR-RV & QLIKE & N & Oxford \\
Luong \& Dok. & 2021 & SPX & 2000--2019 & RF & HAR-RV & MSE & N & Oxford \\
Taylor (2019) & 2019 & Multi & 1990--2017 & NN-CAViaR & CAViaR & VaR cov. & N & Public \\
Cont et al. & 2022 & Multi & 2005--2020 & cGAN & Hist.\ Sim & VaR cov. & N & Various \\
Bussmann et al. & 2021 & Credit & 2005--2019 & GBDT+SHAP & Logistic & AUC & Y & Private \\
Holopainen (2017) & 2017 & Banking & 1970--2014 & RF & Logistic & AUC & N & BIS \\
Giudici et al. & 2022 & Banks & 2008--2020 & GNN & CoVaR & MES & N & Private

\end{longtable}


\end{document}
